% TOP_PROBLEM: {"id": 20001090, "title": "A diameter-free Property B bound for arbitrary van der Waerden sets", "category_id": 2, "category": {"id": 2, "name": "combinatorics", "display_name": "Combinatorics", "description": "Counting problems, graph theory, discrete structures.", "slug": "combinatorics", "order_index": 2, "created_at": "2026-07-31T15:26:25.671Z"}, "status": "partially_solved", "problem_number": "AIM-COMBINATORICS-0215", "set_id": 14}
% FIRST_POSTED: 2026-10-01
% ENTRY_KIND: statement_only
% UnsolvedMath ID 20001090; source catalogue code AIM-COMBINATORICS-0215
% !TeX program = lualatex
% Definition and sources only; this file is not an LLM solution attempt.
\documentclass[11pt,a4paper]{article}
\usepackage[margin=25mm]{geometry}
\usepackage{fontspec}
\setmainfont{lmroman10-regular.otf}[BoldFont=lmroman10-bold.otf,
 ItalicFont=lmroman10-italic.otf,BoldItalicFont=lmroman10-bolditalic.otf]
\usepackage{amsmath,amssymb,mathtools}
\usepackage{unicode-math}
\setmathfont{latinmodern-math.otf}
\AtBeginDocument{\let\setminus\smallsetminus}
\usepackage{xurl,hyperref}
\hypersetup{colorlinks=true,urlcolor=blue}
\setcounter{secnumdepth}{0}
\setlength{\parindent}{0pt}
\setlength{\parskip}{5pt}
\setlength{\emergencystretch}{3em}
\begin{document}
\section*{A diameter-free Property B bound for arbitrary van der Waerden sets}\label{20001090}
{\small UnsolvedMath ID 20001090 · AIM-COMBINATORICS-0215 · Combinatorics · AIM Workshop Problem Lists}
\subsection{Definitions and mathematical statement}
For $k\ge2$, a $k$-term arithmetic progression means
$\{a,a+d,\ldots,a+(k-1)d\}\subseteq\mathbb Z$ with $d>0$.
A finite set $X\subseteq\mathbb Z$ is two-color $k$-progression
forcing if every coloring $X\to\{\mathrm{red},\mathrm{blue}\}$
contains a monochromatic progression of this form entirely within $X$.
Let
\[
 W^*(k)=\min\{|X|:X\subseteq\mathbb Z\text{ is finite and
                                      two-color }k\text{-progression forcing}\}.
\]
Also let $W(k)$ be the minimum $N$ for which $[N]$ is two-color
$k$-progression forcing. These numbers are finite, and
$W^*(k)\le W(k)$ because intervals are among the permitted sets.
The extremal set for $W^*(k)$ need not be an interval and has no
prescribed diameter or maximum element.

The two questions are whether the nonnegative difference
$W(k)-W^*(k)$ is unbounded as $k\to\infty$, and whether
\[
                         \lim_{k\to\infty}\frac{W^*(k)}{W(k)}=1.
\]
Unbounded difference means that no fixed constant bounds the gaps
for all $k$; it does not assert that the gap increases monotonically
or tends to infinity along every sequence. The two proposed behaviours
can hold together, since an unbounded additive saving may still be
a vanishing relative saving. Both questions compare the smallest
cardinality of any forcing set with the forcing threshold for intervals.
\subsection{Short English statement}
How much can arbitrary integer sets save over intervals in forcing monochromatic progressions: is the additive saving unbounded, and does the relative saving vanish?
\subsection{Sources}
\begin{itemize}
\item[S1] ULAM.AI and the UnsolvedMath Contributors, supplied problems.json, record 20001090 (AIM-COMBINATORICS-0215), AIM Workshop Problem Lists. Catalogue identity and source transcription; CC BY 4.0.
\url{https://huggingface.co/datasets/ulamai/UnsolvedMath}.
\item[S2] American Institute of Mathematics, Recent trends in additive combinatorics, item 1.9. Original workshop question underlying AIM-COMBINATORICS-0215.
\url{https://aimath.org/WWN/additivecomb/additivecomb.pdf}.
\end{itemize}
\end{document}
